% ---------------------------------------------------------
\section{Bayesian Influence Functions}\label{sec:bayesian-if}

Classical influence functions give a first-order approximation to the change in a \emph{point estimate} $w^*$ when the training data is perturbed. The Bayesian reformulation replaces this point estimate by a \emph{distribution} over parameters: if we perturb the training data, how does our posterior over the parameter space change? This leads naturally to the Bayesian generalization of influence functions introduced in \citet{kreer2025bif}. In the setting of \Cref{eq:holy-maps} we set $\pi(\beta)$ not to be a single weight, but a whole distribution over parameters $p_{\beta}$, namely the posterior.

\paragraph{Susceptibilities.} The Bayesian influence function is one instance of a more general object studied by \citet{baker2025studying}. Given a tempered Gibbs posterior $p_\beta \propto e^{-\beta L(w)}\varphi(w)$ at inverse temperature $\beta$, any pair of observables $\phi,\psi:W\to\R$ defines a \emph{susceptibility}
\begin{equation}\label{eq:susceptibility}
  \chi(\phi\,|\,\psi) \;:=\; \frac{\partial}{\partial h}\,\E_{p_h}[\phi(w)]\bigg|_{h=0}, \qquad p_h \;\propto\; e^{-\beta\bigl(L(w) - h\,\psi(w)\bigr)}\,\varphi(w).
\end{equation}
The differentiation argument we will apply in \Cref{ex:bif-covariance} shows $\chi(\phi\,|\,\psi) = \beta\,\mathrm{Cov}_{p_\beta}(\phi,\psi)$ in full generality. Taking $\psi = -L(\cdot,z_i)$ gives the Bayesian influence function developed below. Other choices of $(\phi,\psi)$ yield other sensitivity notions, see e.g. the refined learning coefficients \citep[Appendix~D.2]{baker2025studying}. We will not need the general framework again; we mention it only to locate the BIF within the wider susceptibility family.

\subsection{Bayesian influence functions}\label{subsec:bif}

We now specialize the susceptibility story to data attribution. Instead of asking ``how does $\phi(w^*)$ change when we perturb the training weights?'' (\Cref{def:if}) we ask ``how does $\E_{w\sim p}[\phi(w)]$ change?'' The answer turns out to be remarkably clean: the derivative becomes a covariance, and the Hessian inversion of the classical formula is replaced by posterior sampling.

\paragraph{The tempered posterior.} Given training data $\mathcal{D} = \{z_1,\dots,z_n\}$ with per-sample losses $L(w,z_i)$, sample weights $\beta = (\beta_1,\dots,\beta_n)$, and a prior $\varphi(w)$, define the \emph{tempered posterior}
\begin{equation}\label{eq:tempered-posterior}
  p_\beta(w\mid\mathcal{D}) \;\propto\; \exp\!\Bigl(-\sum_{i=1}^n \beta_i\, L(w,z_i)\Bigr)\,\varphi(w).
\end{equation}
When $\beta = \mathbf{1}$ and the loss is a negative log-likelihood, this is the standard Bayesian posterior. When the losses are not log-likelihoods, $p_\beta$ is a \emph{Gibbs measure}; the mathematics is identical.

\begin{definition}[Bayesian influence function]\label{def:bif}
The \emph{Bayesian influence function} (BIF) of training example $z_i$ on the observable $\phi$ is the derivative of the posterior expectation of $\phi$ with respect to the sample weight $\beta_i$:
\[
  \mathrm{BIF}(z_i,\phi) \;:=\; \frac{\partial}{\partial\beta_i}\,\E_{w\sim p_\beta}[\phi(w)]\bigg\rvert_{\beta=\mathbf{1}}.
\]
\end{definition}

The definition is the natural Bayesian analogue of \Cref{def:if}: replace $\phi(w^*(\beta))$ with $\E_{p_\beta}[\phi(w)]$. The following exercise shows that this derivative has a clean closed form. Recall that for two real-valued random variables $X$ and $Y$ defined on the same probability space, the \emph{covariance} is
\[
  \mathrm{Cov}(X,Y) \;=\; \E[XY] - \E[X]\,\E[Y].
\]
When $X$ and $Y$ are functions of the random parameter $w\sim p$, we write $\mathrm{Cov}_{w\sim p}(X(w),Y(w))$ to indicate that the expectation is taken over the posterior $p$. If $\phi$ is vector-valued ($\phi:W\to\R^d$), the covariance $\mathrm{Cov}(X,\phi)$ is the vector whose $j$-th entry is $\mathrm{Cov}(X,\phi_j)$.

\begin{exercise}[The BIF is a covariance]\label{ex:bif-covariance}
Let $p_\beta(w\mid\mathcal{D})$ be the tempered posterior of \Cref{eq:tempered-posterior} and write $p = p_{\mathbf{1}}$ for the unperturbed posterior. By differentiating $\E_{p_\beta}[\phi(w)] = \int \phi(w)\,p_\beta(w\mid\mathcal{D})\,dw$ with respect to $\beta_i$, show that
\begin{equation}\label{eq:bif-covariance}
  \boxed{\;\mathrm{BIF}(z_i,\phi) \;=\; -\,\mathrm{Cov}_{w\sim p}\bigl(L(w,z_i),\;\phi(w)\bigr).\;}
\end{equation}

\emph{Hint:} Write $p_\beta \propto e^{-\sum_j \beta_j L(w,z_j)}\varphi(w)$ and differentiate the ratio $\int \phi\, p_\beta\,dw \,/\, \int p_\beta\,dw$ using the quotient rule. The key identity is $\partial_{\beta_i}\log Z(\beta) = -\E_{p_\beta}[L(w,z_i)]$.
\end{exercise}

\begin{remark}[Statistical physics]\label{rem:stat-phys}
The identity $\partial_{\beta_i}\E[\phi] = -\mathrm{Cov}(L_i,\phi)$ is a standard fluctuation--response relation in statistical physics: the response of an observable to a change in an external field equals the covariance of that observable with the conjugate energy.
\end{remark}

\paragraph{Localization.} Computing expectations over the global posterior $p(w\mid\mathcal{D})$ is generally intractable for neural networks, and we are typically interested in the behavior of a specific trained checkpoint $w^*$, not a global Bayesian average. Following \citet{kreer2025bif}, we \emph{localize} the posterior by adding a Gaussian penalty centered at $w^*$:
\begin{equation}\label{eq:local-posterior}
  p_\gamma(w\mid\mathcal{D},w^*) \;\propto\; \exp\!\Bigl(-\sum_{i=1}^n L(w,z_i) - \frac{\gamma}{2}\|w - w^*\|^2\Bigr),
\end{equation}
where $\gamma > 0$ is a \emph{localization strength} that controls how tightly the posterior concentrates around $w^*$. The \emph{local BIF} is then
\begin{equation}\label{eq:local-bif}
  \mathrm{BIF}_\gamma(z_i,\phi) \;=\; -\,\mathrm{Cov}_\gamma\bigl(L(w,z_i),\;\phi(w)\bigr),
\end{equation}
where $\mathrm{Cov}_\gamma$ denotes covariance under $p_\gamma$. We will see in \Cref{ex:bif-expansion} that the localization parameter $\gamma$ plays exactly the role of the damping parameter $\lambda$ in the PBRF formula \eqref{eq:if-pbrf}.

\paragraph{Practical computation via SGLD.} The local BIF requires estimating a covariance under $p_\gamma$. This is done by drawing samples from $p_\gamma$ using stochastic gradient Langevin dynamics (SGLD), which adds calibrated Gaussian noise to SGD updates:
\[
  w_{t+1} \;=\; w_t - \frac{\epsilon}{2}\Bigl(\frac{n}{m}\sum_{k\in B_t}\nabla_w L(w_t,z_k) + \gamma(w_t - w^*)\Bigr) + \eta_t, \qquad \eta_t\sim\mathcal{N}(0,\epsilon\, I),
\]
where $B_t$ is a mini-batch of size $m$ and $\epsilon$ is the step size. After a burn-in period, the iterates $\{w_t\}$ are approximately distributed according to $p_\gamma$, and the covariance in \eqref{eq:local-bif} is estimated by the sample covariance of $(L(w_t,z_i),\phi(w_t))$ across draws. Running multiple independent chains from $w^*$ improves coverage.

\begin{example}[Bayesian linear regression]\label{ex:bayesian-ols}
Take the linear regression setup of \Cref{ex:ols-if} with a Gaussian prior $w\sim\mathcal{N}(0,\tau^2I)$ added. The posterior is conjugate: $w\mid\mathcal{D}\sim\mathcal{N}(\mu,V)$ with $V = (\tau^{-2}I + \sigma^{-2}A)^{-1}$ and $\mu = \sigma^{-2}V\sum_i x_i y_i$. A direct Gaussian-moment computation applied to \eqref{eq:bif-covariance} with the observable $\phi(w) = w$ gives
\[
  \mathrm{BIF}(z_j, w) \;=\; -\mathrm{Cov}_p\bigl(L(w,z_j),\,w\bigr) \;=\; \frac{V\, x_j\, r_j}{\sigma^2} \;=\; (A + \lambda I)^{-1}\,x_j\,r_j, \qquad \lambda := \sigma^2/\tau^2,
\]
where $r_j = y_j - x_j^\top\mu$ is the posterior-mean residual. This is the \emph{damped} influence function with the damping parameter $\lambda$ equal to the ratio of observation noise to prior variance. In the flat-prior limit $\tau\to\infty$ we have $\lambda\to 0$, $\mu\to w^*$, and we recover the classical IF of \Cref{ex:ols-if}.
\end{example}

\subsection{Exercise: Connecting Bayesian and classical influence functions}\label{subsec:bif-if-connection}
%TODO: Specify this is the non degenerate case
As we have seen in the SLT day, tools involving the Hessian are usually a low-order approximation of an expansion. This is also true for influence functions: In the non-degenerate case the classical IF is the leading-order term of the BIF under a Laplace approximation. The following exercise makes this precise by computing a power series expansion of the BIF covariance and identifying the classical IF as the leading-order term. The argument follows \citet[Appendix~A]{kreer2025bif}.

\begin{exercise}[Power series expansion of Bayesian influence functions]\label{ex:bif-expansion}
Let $w^*$ be a local minimum of the training loss $L(w) = \sum_{i=1}^n L(w,z_i)$ with positive-definite Hessian $H = \nabla_w^2 L(w^*)$. Write $\Delta w = w - w^*$, and abbreviate $g_\phi = \nabla_w\phi(w^*)$, $H_\phi = \nabla_w^2\phi(w^*)$ for the gradient and Hessian of the observable, and $g_i = \nabla_w L(w^*,z_i)$, $H_i = \nabla_w^2 L(w^*,z_i)$ for the per-sample loss.
\begin{enumerate}
  \item[(a)] \textbf{(Laplace approximation.)} Consider the posterior $p(w\mid\mathcal{D}) \propto e^{-L(w)}\,\varphi(w)$. Taylor-expand $L(w)$ to second order around $w^*$, using the fact that $\nabla L(w^*) = 0$ at a minimum. Argue that for large $n$ the quadratic term in $L$ dominates the prior $\varphi$, and conclude that
  \[
    p(w\mid\mathcal{D}) \;\approx\; \mathcal{N}(w^*,\;H^{-1}).
  \]
  Under this approximation, $\Delta w \sim \mathcal{N}(0, H^{-1})$. (This is called the \emph{Laplace approximation}; the Bernstein--von Mises theorem guarantees it is asymptotically exact for non-singular models.)

  \item[(b)] \textbf{(Taylor expansion of the covariance.)} Expand $\phi(w)$ and $L(w,z_i)$ in Taylor series around $w^*$:
  \begin{align*}
    \phi(w) &= \phi(w^*) + g_\phi^\top\Delta w + \tfrac{1}{2}\Delta w^\top H_\phi\,\Delta w + \cdots, \\
    L(w,z_i) &= L(w^*,z_i) + g_i^\top\Delta w + \tfrac{1}{2}\Delta w^\top H_i\,\Delta w + \cdots.
  \end{align*}
  Using the fact that constant terms drop out of covariances and that $\mathrm{Cov}(X,Y) = \sum_{k,m\ge 1}\mathrm{Cov}(T_k[\phi],\,T_m[L_i])$ where $T_k[f]$ denotes the $k$-th order term in the Taylor expansion of $f$, show that the BIF decomposes as
  \[
    \mathrm{BIF}(z_i,\phi) = -\sum_{\substack{k,m\ge 1\\k+m\text{ even}}} \mathrm{Cov}_{\mathcal{N}}\bigl(T_k[\phi],\;T_m[L_i]\bigr).
  \]
  Why do terms with $k+m$ odd vanish?

  \item[(c)] \textbf{(Leading order: classical IF.)} Compute the $(k,m) = (1,1)$ term:
  \[
    \mathrm{Cov}_{\mathcal{N}}\bigl(g_\phi^\top\Delta w,\;g_i^\top\Delta w\bigr) \;=\; g_\phi^\top H^{-1} g_i.
  \]
  Conclude that $-g_\phi^\top H^{-1}g_i = -\nabla_w\phi(w^*)^\top H^{-1}\nabla_w L(w^*,z_i)$ is the classical influence function $\mathcal{I}(z_i,\phi)$ from \Cref{def:if}. This is the leading-order term of the BIF.

  \item[(d)] \textbf{(Second-order correction.)} Compute the $(k,m) = (2,2)$ term using Isserlis' theorem (the Gaussian moment identity $\E[\Delta w_a\Delta w_b\Delta w_c\Delta w_d] = \Sigma_{ab}\Sigma_{cd} + \Sigma_{ac}\Sigma_{bd} + \Sigma_{ad}\Sigma_{bc}$ with $\Sigma = H^{-1}$). Show that
  \[
    \mathrm{Cov}_{\mathcal{N}}\!\left(\tfrac{1}{2}\Delta w^\top H_\phi\,\Delta w,\;\tfrac{1}{2}\Delta w^\top H_i\,\Delta w\right) = \tfrac{1}{2}\,\tr\bigl(H_\phi\,H^{-1}\,H_i\,H^{-1}\bigr).
  \]
  This correction involves the \emph{Hessians} of the observable and per-sample loss. It captures second-order curvature interactions that the classical IF misses entirely. In the linear regression setting of \Cref{ex:bayesian-ols}, why does this correction vanish?

  \item[(e)] \textbf{(Localized version: damped IF.)} Now consider the local BIF of \Cref{eq:local-bif} with localization strength $\gamma$. The localized posterior is approximately $\mathcal{N}(w^*,\,(H+\gamma I)^{-1})$. Repeat the leading-order computation of part (c) to show that
  \[
    \mathrm{BIF}_\gamma(z_i,\phi) \;\approx\; -\,\nabla_w\phi(w^*)^\top\,(H + \gamma I)^{-1}\,\nabla_w L(w^*,z_i).
  \]
  This is precisely the damped influence function of \Cref{eq:if-pbrf}, with the localization strength $\gamma$ playing the role of the damping parameter $\lambda$. The local BIF is thus a natural, higher-order generalization of the damped IF: it agrees at leading order and includes all the corrections from parts (b)--(d) with $H^{-1}$ replaced by $(H + \gamma I)^{-1}$.

\end{enumerate}
\end{exercise}
\begin{remark}
    We conclude that the classical (damped) IF is the leading term of the (local) BIF under a Laplace approximation. For non-singular models with enough data, the Laplace approximation converges and the higher-order terms are small. But for singular models (neural networks), the Laplace approximation fails, the posterior is not Gaussian, and the BIF captures geometry that the classical IF cannot see.
\end{remark}
\subsection{Long Exercise: Influence functions as optimal linear transport}\label{subsec:bif-ot}

The previous exercise showed that the classical IF is the leading-order term of the BIF under a Laplace approximation. We now take a different perspective, due to \citet{mlodozeniec2025distributional}: instead of asking ``what does the BIF reduce to?'' we ask ``what is the \emph{best} way to shift the parameters to approximate the effect of a data perturbation?'' The answer turns out to be: the influence function, and its optimality does not require the Laplace approximation to hold.

\paragraph{Setup.} Let $L(w) = \sum_{i=1}^n L(w,z_i)$ denote the total training loss, and abbreviate $L_k(w) := L(w,z_k)$ for the per-sample loss at example $k$. Define the \emph{Boltzmann distributions} at inverse temperature $\beta$:
\begin{align*}
  p_0^\beta(w) &\propto e^{-\beta\,L(w)}, &
  p_\varepsilon^\beta(w) &\propto e^{-\beta\bigl(L(w)\,-\,\varepsilon\,L_k(w)\bigr)},
\end{align*}
so that $p_\varepsilon^\beta$ is the Boltzmann distribution after downweighting sample $k$ by~$\varepsilon$. The low-temperature limit $\beta\to\infty$ concentrates $p_0^\beta$ on the set of global minima $\mathcal{S}_L = \{w : L(w) = \inf L\}$.

Now suppose we want to approximate the effect of the perturbation without resampling from $p_\varepsilon^\beta$. The simplest approach is to take each parameter $w$ drawn from $p_0^\beta$ and \emph{shift} it to $w + \varepsilon\,r(w)$ for some smooth vector field $r:\R^d\to\R^d$. If $r$ is chosen well, the distribution of the shifted parameters should be close to $p_\varepsilon^\beta$. The question is: what is the optimal $r$?

To make this precise, we need to know the density of the shifted parameters. If $w\sim p_0^\beta$ and we set $\tilde w = w + \varepsilon\,r(w)$, then by the multivariate change-of-variables formula, $\tilde w$ has density
\begin{equation}\label{eq:shifted-density}
  q_\varepsilon(\tilde w) \;=\; \frac{p_0^\beta\bigl(\tilde w - \varepsilon\,r(\tilde w) + O(\varepsilon^2)\bigr)}{\bigl|\det\bigl(I + \varepsilon\,\nabla r(w)\bigr)\bigr|}\,,
\end{equation}
where $\nabla r$ is the $d\times d$ Jacobian matrix of $r$ and the denominator is the usual volume-change factor.\footnote{For small $\varepsilon$ the map $w\mapsto w+\varepsilon r(w)$ is a diffeomorphism, so the inverse is well-defined. The exact formula requires evaluating $r$ at the pre-image; the expression above is correct to the order we need.} We measure how well $q_\varepsilon$ approximates $p_\varepsilon^\beta$ by the KL divergence $D_{\mathrm{KL}}(q_\varepsilon\|p_\varepsilon^\beta)$. Rather than computing this directly, it is cleaner to write it as an expectation over the \emph{original} distribution $p_0^\beta$:
\begin{equation}\label{eq:kl-expectation}
  D_{\mathrm{KL}}(q_\varepsilon\|p_\varepsilon^\beta) \;=\; \E_{w\sim p_0^\beta}\!\Bigl[\log p_0^\beta(w) - \log p_\varepsilon^\beta\bigl(w + \varepsilon r(w)\bigr) - \log\bigl|\det\bigl(I + \varepsilon\nabla r(w)\bigr)\bigr|\Bigr].
\end{equation}
This identity follows by substituting \eqref{eq:shifted-density} into the definition of KL and changing variables back to $w$. We define the \emph{asymptotic KL cost}
\begin{equation}\label{eq:ot-cost}
  \mathcal{F}(r,\varepsilon) \;:=\; \lim_{\beta\to\infty}\,\frac{1}{\beta}\,D_{\mathrm{KL}}(q_\varepsilon\|p_\varepsilon^\beta).
\end{equation}
The $1/\beta$ normalization kills the determinant term (which is $O(1)$) and the entropy terms, isolating the energy contribution (which is $O(\beta)$).

\begin{exercise}[Influence functions as optimal parameter shifts]\label{ex:bif-ot}
Let $w^*$ be a non-degenerate minimum of $L$ with positive-definite Hessian $H = \nabla^2 L(w^*)$. In this exercise we work at the unique minimum, so the $\beta\to\infty$ limit simply evaluates everything at $w^*$.
\begin{enumerate}
  \item[(a)] \textbf{(KL in terms of energies.)} Starting from \eqref{eq:kl-expectation}, substitute $p_0^\beta(w) = e^{-\beta L(w)}/Z(\beta,0)$ and $p_\varepsilon^\beta(w) = e^{-\beta(L(w)-\varepsilon L_k(w))}/Z(\beta,\varepsilon)$, and show that
  \begin{align*}
    \frac{1}{\beta}\,D_{\mathrm{KL}} \;=\; &\E_{p_0^\beta}\!\bigl[L(w + \varepsilon r) - \varepsilon\,L_k(w + \varepsilon r) - L(w)\bigr] \\
    &+\; \frac{1}{\beta}\log\frac{Z(\beta,\varepsilon)}{Z(\beta,0)} \;-\; \frac{1}{\beta}\,\E_{p_0^\beta}\!\bigl[\log|\det(I + \varepsilon\nabla r)|\bigr],
  \end{align*}
  where $Z(\beta,\varepsilon) = \int e^{-\beta(L(w)-\varepsilon L_k(w))}\,dw$. Explain why the last two terms vanish in the $\beta\to\infty$ limit. \emph{Hint:} $\log|\det(I+\varepsilon\nabla r)| = O(1)$, and by Laplace's method $\tfrac{1}{\beta}\log Z(\beta,\varepsilon) \to -\inf_w[L(w) - \varepsilon L_k(w)]$.

  \item[(b)] \textbf{(Taylor expansion at $w^*$.)} Since the expectation under $p_0^\beta$ concentrates at $w^*$ as $\beta\to\infty$, expand $L(w^* + \varepsilon r)$ and $\varepsilon\,L_k(w^* + \varepsilon r)$ using $\nabla L(w^*) = 0$:
  \begin{align*}
    L(w^* + \varepsilon r) &= L(w^*) + \tfrac{\varepsilon^2}{2}\,r^\top H\,r + O(\varepsilon^3), \\
    \varepsilon\,L_k(w^* + \varepsilon r) &= \varepsilon\,L_k(w^*) + \varepsilon^2\,\nabla L_k(w^*)^\top r + O(\varepsilon^3).
  \end{align*}
  Similarly, show that $\inf_w[L(w) - \varepsilon L_k(w)] = L(w^*) - \varepsilon\,L_k(w^*) - \tfrac{\varepsilon^2}{2}\,\nabla L_k^\top H^{-1}\nabla L_k + O(\varepsilon^3)$, by minimizing the Taylor expansion of $L - \varepsilon L_k$ around $w^*$. Combine everything to obtain
  \begin{equation}\label{eq:ot-quadratic}
    \mathcal{F}(r,\varepsilon) \;=\; \frac{\varepsilon^2}{2}\,\bigl(r - H^{-1}\nabla L_k\bigr)^\top H\,\bigl(r - H^{-1}\nabla L_k\bigr) \;+\; O(\varepsilon^3).
  \end{equation}

  \emph{Hint:} After substitution, the terms that do not depend on $r$ should assemble into $\tfrac{\varepsilon^2}{2}\nabla L_k^\top H^{-1}\nabla L_k$ and cancel with the partition-function contribution, leaving a perfect square.

  \item[(c)] \textbf{(IFs are optimal.)} The cost \eqref{eq:ot-quadratic} is a squared Mahalanobis distance (in the Hessian metric) between the shift $r$ and the influence function $r_{\mathrm{IF}} := H^{-1}\nabla L_k(w^*)$. Conclude that the IF minimizes $\mathcal{F}$ to $O(\varepsilon^2)$ among \emph{all} smooth shift maps $w\mapsto w+\varepsilon r(w)$, and that the minimum cost is $O(\varepsilon^3)$.

  State the result in words: \emph{the influence function is the approximately KL-optimal way to shift the parameters of the unperturbed Boltzmann distribution to approximate the perturbed one.}

  \item[(d)] \textbf{(Degenerate case.)} Now suppose $w^*$ lies on a minimum manifold $\mathcal{S}_L$ and $H$ is singular with null space $N$. Argue (without detailed proof) that:
  \begin{itemize}
    \item The cost \eqref{eq:ot-quadratic} generalizes to $\mathcal{F}(r,\varepsilon) = \tfrac{\varepsilon^2}{2}(r - H^+\nabla L_k)^\top H\,(r - H^+\nabla L_k) + O(\varepsilon^3)$, where $H^+$ is the Moore--Penrose pseudoinverse.
    \item The minimizer is $r = H^+\nabla L_k + r_{\mathrm{NS}}$ for any $r_{\mathrm{NS}}\in N$: directions in the null space of $H$ have zero cost, because the Hessian assigns zero curvature to movement along the minimum manifold.
    \item The IF (with pseudoinverse) is therefore optimal for the directions that matter, but the component of the shift along the flat directions is unconstrained.
  \end{itemize}
  Connect this to the SLT perspective: for singular models, the minimum is a variety, and the IF tells you how to move \emph{off} the variety but is silent about movement \emph{along} it.

  \item[(e)] \textbf{(Linear regression check.)} Specialize to the Bayesian linear regression setting of \Cref{ex:bayesian-ols} with prior $w\sim\mathcal{N}(0,\tau^2I)$. The Boltzmann distribution at inverse temperature $\beta = 1$ is the Bayesian posterior $\mathcal{N}(\mu,V)$, which is already Gaussian, no need to take $\beta\to\infty$. A constant shift $w\mapsto w + \varepsilon r$ sends $\mathcal{N}(\mu,V)$ to $\mathcal{N}(\mu + \varepsilon r,\,V)$. The true perturbed posterior (downweighting $z_k$ by $\varepsilon$) has mean $\mu + \varepsilon\,\mathrm{BIF}(z_k,w) + O(\varepsilon^2)$ and covariance $V + O(\varepsilon)$. Using the KL formula for Gaussians with the same covariance, show that
  \[
    D_{\mathrm{KL}}\bigl(\mathcal{N}(\mu + \varepsilon r,\,V)\;\big\|\;\mathcal{N}(\mu + \varepsilon b,\,V)\bigr) \;=\; \frac{\varepsilon^2}{2}\,(r-b)^\top V^{-1}(r-b),
  \]
  which is minimized at $r = b = \mathrm{BIF}(z_k,w) = (A + \lambda I)^{-1}x_k r_k$. Verify that this is consistent with \eqref{eq:ot-quadratic}: the ``Hessian of the energy'' (training loss plus prior) is $V^{-1} = \sigma^{-2}A + \tau^{-2}I$, and $V^{-1}\cdot\mathrm{BIF} = \sigma^{-2}x_k r_k = \nabla_w L(w,z_k)\big|_{w=\mu}$.

  \item[(f)] \textbf{(Comparison with \Cref{ex:bif-expansion}.)} The Laplace expansion exercise and this exercise both relate the IF to the BIF, but they answer different questions:
  \begin{itemize}
    \item \Cref{ex:bif-expansion} is an \emph{algebraic identity}: BIF $=$ IF $+$ higher-order corrections. It requires the Laplace approximation (hence non-singular $H$, Gaussian posterior) and tells you what happens when you truncate the BIF.
    \item This exercise is an \emph{optimality result}: the IF shift minimizes the KL divergence between the shifted and true perturbed distributions. The $\beta\to\infty$ result (parts a--d) needs only mild regularity of $\mathcal{L}$ and works even for singular $H$ (via the pseudoinverse).
  \end{itemize}
  Summarize: the Laplace expansion tells you \emph{what} the IF is (the leading Gaussian term of the BIF). The optimality perspective tells you \emph{why} it works (it is the best first-order correction to the parameter distribution). The second viewpoint does not require the posterior to be Gaussian, which is why \citet{mlodozeniec2025distributional} argue that it provides a better explanation for the empirical success of IFs in deep learning.
\end{enumerate}
\end{exercise}

\begin{remark}[Optimal transport]\label{rem:optimal-transport}
In the language of optimal transport, the map $w\mapsto w+\varepsilon r(w)$ is called a \emph{transport map}, and the distribution of the shifted parameters is its \emph{pushforward}, written $T_{\varepsilon\#}p_0^\beta$. The cost functional $\mathcal{F}$ is the asymptotic KL divergence between the pushforward and the target. \Cref{ex:bif-ot} then says that the influence function is the \emph{KL-optimal transport map} from the unperturbed to the perturbed Boltzmann distribution, at first order in $\varepsilon$. This is the perspective developed in \citet{mlodozeniec2025distributional}, who prove a more general version (their Theorem~3) that applies to all smooth diffeomorphisms, not just maps of the form $w+\varepsilon r(w)$.
\end{remark}
